\newcommand*{\version}{0.1.0}

\documentclass[12pt,
              article,
              oneside
              ]{memoir}

\usepackage{iftex}

\ifluatex
  \usepackage{fontspec}
\fi

\usepackage[
  % latin,
  % english,
  spanish,% <- principal
  es-minimal,
  es-noshorthands,
  es-tabla,
]{babel}

\ifluatex
% LuaLaTeX --------------------
\defaultfontfeatures{Ligatures=TeX, Scale=MatchLowercase}
\usepackage[MainFeatures={medium,
  Ligatures=Common,
  Numbers=OldStyle,
  SmallCapsFeatures = {LetterSpace=5.0}}]
{junicodevf}
%
% \setsansfont{Libertinus Sans}
% (sin las ligaduras de TeX: -{}-help no se vuelve una raya)
\setmonofont[Scale=MatchLowercase, Ligatures=ResetAll]{DejaVu Sans Mono}
% \setmonofont[Scale=MatchLowercase]{Libertinus Mono}
% \setmonofont[Scale=MatchLowercase}]{Fira Code}
\else
% PDFLaTeX ---------------------
\usepackage[T1]{fontenc}
\usepackage[osf]{libertinus}
\usepackage[scaled=0.9]{FiraMono}
% \usepackage[scaled=0.9]{cascadia-code}
% \usepackage[scaled=0.9]{roboto}
\input glyphtounicode
\pdfgentounicode=1
\fi

\usepackage[babel=true, verbose=false, tracking=true, expansion=true, protrusion=true, final, draft=false]{microtype}
\SetTracking{ encoding = *, shape = sc }{ 30 }
\SetTracking[context = notracking, ]{encoding = *}{0}

\usepackage[autostyle=false, style=british]{csquotes}
\usepackage{latexcolors}

\newcommand{\Rojo}{\color[HTML]{8B0000}}
\newcommand*{\rojo}[1]{\textcolor[HTML]{8B0000}{#1}}
\newcommand*{\rojoit}[1]{\textit{\textcolor[HTML]{8B0000}{#1}}}
\newcommand{\Azul}{\color{bluenivaca}}
\newcommand*{\azul}[1]{\textcolor{bluenivaca}{#1}}
\newcommand{\Verde}{\color{darkgreen}}
\newcommand*{\verde}[1]{\textcolor{darkgreen}{#1}}
\newcommand{\Amarillo}{\color{goldenpoppy}}
\newcommand*{\amarillo}[1]{\textcolor{goldenpoppy}{#1}}
\newcommand{\Morado}{\color{byzantium}}
\newcommand*{\morado}[1]{\textcolor{byzantium}{#1}}
\newcommand{\Naranja}{\color{orange}}
\newcommand*{\naranja}[1]{\textcolor{orange}{#1}}

% \usepackage[margin=3cm]{geometry}

\usepackage{graphicx}
\usepackage{booktabs}

\clubpenalty=9999
\widowpenalty=9999
\brokenpenalty=5000
\frenchspacing
\raggedbottom
\midsloppy
\parindent=14pt
\fboxsep=2pt

\renewcommand\thesection{\arabic{section}}
\chapterstyle{reparticle}

\usepackage{enumitem}
\setlist{noitemsep, nosep}
\setlist[description]{font=\itshape}

\newcommand*{\elipsis}{\,.\,.\,.\,}
\newcommand*{\elipsisb}{\textup{[\elipsis]}}

\renewenvironment{quote}%
 {\begin{list}{}%
              {\setlength\rightmargin{0pt}%
              \setlength\leftmargin{\parindent}}%
  \item[]\small\ignorespaces}
 {\unskip\end{list}}


\definecolor{Link}{HTML}{408080}
\usepackage{hyperref}
\hypersetup{
  colorlinks=true,
  linkcolor=blue,
  filecolor=blue,
  citecolor = black,
  urlcolor=Link,
}
\urlstyle{tt}
\usepackage{bookmark}

\setlength{\droptitle}{-2cm}

\title{\bfseries%
  \texttt{py-tuner} v.\version{} \\
  Manual de usuario}

\author{%
  Nicolás Vaughan\\
  \begin{minipage}[t]{0.8\linewidth}
    \centering\small%
    \href{mailto:nicolasvaughan@proton.me}{\texttt{nicolasvaughan@proton.me}}\\
    \url{https://nicolasvaughan.org}\\
    \url{https://github.com/nivaca/py-tuner}\\
    \url{https://codeberg.org/nivaca/py-tuner}\\
  \end{minipage}}


\date{2026-11-27}

% The pictures, made by tools/manual_images.py
\graphicspath{{img/manual/}{img/}}
\usepackage{longtable}

% the title and the author in the PDF's own metadata
\hypersetup{pdftitle={py-tuner: manual de usuario}, pdfauthor={Nicolás Vaughan}}

% the index (makeindex, with the style manual.ist: see latexmkrc)
\makeindex

% menus, keys, and the pictures of the window
\newcommand*{\menu}[1]{\textit{#1}}
\newcommand*{\then}{\,>\,}
\newcommand*{\key}[1]{\texttt{#1}}
\newcommand*{\shot}[3]{%
  \begin{figure}[htbp]
    \centering
    \includegraphics[width=#2\textwidth,height=0.75\textheight,keepaspectratio]{#1}
    \caption{#3}
    \label{fig:#1}
  \end{figure}}

\begin{document}
\maketitle

\tableofcontents*

\section{Introducción}
\noindent
\texttt{py-tuner} es un afinador para guitarra, bajo, contrabajo y violín,
con un modo cromático para cualquier otro instrumento. Escucha la tarjeta
de sonido, reconoce la nota que suena, la compara con la cuerda más cercana
de la afinación elegida (o con la cuerda que se fije) y muestra cuántos
cents se aparta. Un cent es la centésima parte de un semitono.

\index{precisión}\index{inarmonicidad}%
Lo hice para que fuera preciso. En una cuerda estable, cada cuadro de
análisis cae a unas centésimas de cent de la altura verdadera, y la lectura
que se muestra promedia medio segundo de cuadros. La altura sale de todos
los armónicos de la cuerda a la vez, teniendo en cuenta que una cuerda real
es inarmónica (sus parciales altos suenan un poco agudos); por eso funciona
también cuando un micrófono pequeño apenas recoge la fundamental de una
cuerda grave, como pasa con el bajo y el contrabajo.

\index{interfaces}%
Muestra la desviación con una aguja y con un estroboscopio, y funciona de
tres maneras: en una ventana, en la terminal a pantalla completa y en la
línea de órdenes. La interfaz está en español o en inglés. Este manual
describe la versión \version{}; las imágenes son de la ventana en español.

\section{Instalación}
\label{sec:installing}

\index{instalación}%
Las versiones están en GitHub y en Codeberg (las direcciones están en la
primera página). Cada una trae los archivos que siguen, todos firmados:
véase la sección~\ref{sec:signatures}.

\subsection{Fedora}

\index{fedora@Fedora}\index{copr@Copr}\index{rpm@RPM}%
El RPM está también en Copr, para que \texttt{dnf upgrade} traiga cada
versión nueva:

\begin{quote}
\texttt{sudo dnf copr enable nivaca/py-tuner}\\
\texttt{sudo dnf install py-tuner}
\end{quote}

\noindent Usa NumPy, PortAudio y PySide6 de Fedora, e instala una entrada
en el menú, las páginas de manual (\texttt{man py-tuner}, en español y en
inglés) y este manual.

\subsection{Cualquier Linux: la AppImage}

\index{appimage@AppImage}%
\texttt{py-tuner-\version-x86\_64.AppImage} lleva su propio Python, Qt y
PortAudio, y no hay que instalarla. Funciona en cualquier Linux con glibc
2.34 o posterior (Ubuntu 22.04, Debian 12, Fedora 35 y posteriores):

\begin{quote}
\texttt{chmod +x py-tuner-\version-x86\_64.AppImage}\\
\texttt{./py-tuner-\version-x86\_64.AppImage}\\
\texttt{./py-tuner-\version-x86\_64.AppImage tui}
\end{quote}

\noindent La primera orden abre la ventana; la segunda, la interfaz de
terminal. Cualquier otra orden u opción de la sección~\ref{sec:cli} se da
de la misma manera.

\subsection{Desde el código}

\index{código}\index{python@Python}%
Con Python 3.14 o posterior y la PortAudio del sistema (\texttt{sudo dnf
install portaudio} en Fedora, \texttt{sudo apt install libportaudio2} en
Debian y Ubuntu), el wheel de una versión se instala con \texttt{pip
install py\_tuner-\version-py3-none-any.whl}. Desde una copia del código,
con uv (\url{https://docs.astral.sh/uv/}), \texttt{uv run py-tuner} lo
corre enseguida.

\subsection{Comprobar los archivos de una versión}
\label{sec:signatures}

\index{firmas}%
Cada archivo de una versión tiene su firma (\texttt{ARCHIVO.asc}), y
\texttt{SHA256SUMS} lista las sumas de todos, con su propia firma. Se
comprueban con la clave pública que se publica con la versión:

\begin{quote}
\texttt{gpg -{}-import RPM-GPG-KEY-py-tuner}\\
\texttt{gpg -{}-verify ARCHIVO.asc ARCHIVO}
\end{quote}

\noindent El RPM va firmado también por dentro; \texttt{dnf} lo comprueba
una vez importada la clave (\texttt{sudo rpm -{}-import
RPM-GPG-KEY-py-tuner}).

\section{Cómo abrir el programa}

\index{abrir el programa}%
Desde el menú, \menu{py-tuner} abre la ventana. Desde una terminal:

\begin{quote}
\begin{tabular}{@{}ll@{}}
\texttt{py-tuner} & la ventana\\
\texttt{py-tuner tui} & la interfaz de terminal (sección~\ref{sec:tui})\\
\texttt{py-tuner cli} & una línea de texto (sección~\ref{sec:cli})\\
\texttt{py-tuner tone E2} & una nota de referencia (sección~\ref{sec:tone})\\
\texttt{py-tuner devices} & la lista de entradas de sonido\\
\texttt{py-tuner instruments} & la lista de instrumentos y afinaciones\\
\end{tabular}
\end{quote}

\noindent La ventana se abre en el modo cromático
(sección~\ref{sec:chromatic}) hasta que se elija otro instrumento; la
elección se guarda. \texttt{py-tuner -{}-help} lista todas las opciones. Las opciones
cambian la configuración solo en esa sesión, salvo que se añada
\texttt{-{}-save} (sección~\ref{sec:settings}).

\section{La ventana}
\label{sec:window}

\index{ventana}%
\shot{window_es}{0.75}{La ventana: la cuerda Sol de una guitarra, 0,4 cents alta}

\noindent De arriba abajo, la ventana tiene:

\begin{description}
\item[La barra] el instrumento, su afinación (con sus cuerdas), el La de
  referencia (sección~\ref{sec:a4}) y la configuración.
\item[La nota] la nota con que se compara el sonido, en letras grandes;
  debajo, la desviación en cents y la frecuencia que se oye, con una flecha
  hacia la frecuencia que debería tener.
\item[La aguja] la desviación, de $-50$ a $+50$ cents. La estrecha zona
  verde del centro es la tolerancia: una nota dentro de ella está afinada
  (sección~\ref{sec:settings}).
\item[El estroboscopio] tres bandas de franjas que se mueven a la derecha
  cuando la nota está alta y a la izquierda cuando está baja, y se quedan
  quietas cuando está afinada (sección~\ref{sec:strobe}).
\item[La historia] la desviación en los últimos diez segundos, la más
  reciente a la derecha, sobre la banda de la tolerancia
  (sección~\ref{sec:history}).
\item[Las cuerdas] un botón por cada cuerda de la afinación. La cuerda con
  que se compara va enmarcada.
\item[La barra de estado] la entrada en uso, su frecuencia de muestreo y el
  nivel del sonido, con una marca en el nivel por debajo del cual los
  sonidos se ignoran.
\end{description}

\index{colores}%
Los colores dicen lo mismo en todas partes: verde dentro de la tolerancia,
ámbar dentro de cinco veces la tolerancia, rojo más allá.

\shot{flat_es}{0.75}{El Mi grave de un bajo de cinco cuerdas, 12,3 cents bajo}

\index{cuerda fija}%
La cuerda con que se compara es, normalmente, la más cercana al sonido. Un
clic en el botón de una cuerda la fija: desde entonces todo sonido se
compara con esa cuerda, por lejos que esté, lo que ayuda cuando una cuerda
está muy floja y suena más cerca de la siguiente. Otro clic en el mismo
botón la suelta.

\shot{locked_es}{0.75}{La cuerda La del violín fija, con el La a 442~Hz}

\section{Instrumentos y afinaciones}
\label{sec:tunings}

\index{instrumentos}\index{afinaciones}%
Cada instrumento tiene primero su afinación habitual, y algunas otras:

\begin{longtable}{@{}lll@{}}
\toprule
Instrumento & Afinación & Cuerdas, de la más grave\\
\midrule
\endhead
Guitarra & Estándar & \texttt{E2 A2 D3 G3 B3 E4} \\
 & Drop D & \texttt{D2 A2 D3 G3 B3 E4} \\
 & Medio tono abajo & \texttt{Eb2 Ab2 Db3 Gb3 Bb3 Eb4} \\
 & Un tono abajo & \texttt{D2 G2 C3 F3 A3 D4} \\
 & Drop C & \texttt{C2 G2 C3 F3 A3 D4} \\
 & DADGAD & \texttt{D2 A2 D3 G3 A3 D4} \\
 & Abierta en Re & \texttt{D2 A2 D3 F\#3 A3 D4} \\
 & Abierta en Sol & \texttt{D2 G2 D3 G3 B3 D4} \\
 & Abierta en Mi & \texttt{E2 B2 E3 G\#3 B3 E4} \\
 & Abierta en Do & \texttt{C2 G2 C3 G3 C4 E4} \\
 & 7 cuerdas & \texttt{B1 E2 A2 D3 G3 B3 E4} \\
\addlinespace
Bajo & Estándar & \texttt{E1 A1 D2 G2} \\
 & Drop D & \texttt{D1 A1 D2 G2} \\
 & Medio tono abajo & \texttt{Eb1 Ab1 Db2 Gb2} \\
 & 5 cuerdas & \texttt{B0 E1 A1 D2 G2} \\
 & 5 cuerdas con Do agudo & \texttt{E1 A1 D2 G2 C3} \\
 & 6 cuerdas & \texttt{B0 E1 A1 D2 G2 C3} \\
\addlinespace
Contrabajo & Orquestal & \texttt{E1 A1 D2 G2} \\
 & Solista & \texttt{F\#1 B1 E2 A2} \\
 & 5 cuerdas (Si grave) & \texttt{B0 E1 A1 D2 G2} \\
 & 5 cuerdas (Do grave) & \texttt{C1 E1 A1 D2 G2} \\
 & Por quintas & \texttt{C1 G1 D2 A2} \\
\addlinespace
Violín & Estándar & \texttt{G3 D4 A4 E5} \\
 & Cruzada (La Mi La Mi) & \texttt{A3 E4 A4 E5} \\
 & Scordatura Sol Re Sol Re & \texttt{G3 D4 G4 D5} \\
\bottomrule
\end{longtable}

\noindent El instrumento fija también el margen en el que se busca la
fundamental, algo más allá de su cuerda más grave, para que una cuerda muy
floja se siga encontrando.

\index{afinaciones propias}%
\label{sec:own-tunings}
Las afinaciones propias se añaden al archivo de configuración
(sección~\ref{sec:files}), con el programa cerrado, como una lista
\texttt{custom\_tunings}. Cada afinación dice su instrumento
(\texttt{guitar}, \texttt{bass}, \texttt{double-bass} o \texttt{violin}),
su nombre y sus cuerdas, de la más grave a la más aguda:

\begin{verbatim}
"custom_tunings": [
  {"instrument": "guitar", "name": "Nashville",
   "strings": "E3 A3 D4 G4 B3 E4"}
]
\end{verbatim}

\noindent Aparecen después de las afinaciones del instrumento. Una
afinación con una nota que no se puede leer se deja fuera.

\section{El La de referencia}
\label{sec:a4}

\index{La de referencia}\index{a4@A4}%
Todas las notas se miden contra el La que está sobre el Do central (A4),
440~Hz si no se cambia. Puede ir de 380 a 480~Hz, en décimas de hercio: en
la ventana, con la casilla \menu{A4} (las flechas lo mueven 0,1~Hz, Re Pág
y Av Pág 1~Hz) o con el menú de al lado, que tiene los habituales (415,
430, 432, 435, de 438 a 444, 446 y 466~Hz); en la terminal, con las
flechas; en la línea de órdenes, con \texttt{-a}:

\begin{quote}
\texttt{py-tuner -a 442}
\end{quote}

\noindent Subir el La 1~Hz sube cada nota unos 3,9 cents.

\section{El modo cromático y los nombres de las notas}
\label{sec:chromatic}

\index{modo cromático}%
El instrumento \menu{Cromático}, con el que arranca el programa, no tiene
cuerdas: el sonido se compara con la nota más cercana de la escala
temperada, de A0 a G6. Sirve para
cualquier instrumento, y para revisar las notas a lo largo de una cuerda.

\index{nombres de las notas}%
Las notas pueden escribirse en inglés (C, D, E\elipsis), latinos (Do, Re,
Mi\elipsis) o alemanes (con H para Si y B para Si$\flat$), con
sostenidos, bemoles o ambos (F$\sharp$/G$\flat$); la octava es siempre la científica, en la que el Do central es C4
(Do4).

\shot{chromatic_es}{0.75}{El modo cromático, con nombres latinos y bemoles}

\section{La aguja y el estroboscopio}
\label{sec:strobe}

\index{aguja}%
La aguja muestra la desviación de inmediato, y es la más fácil de leer de
lejos. Sigue la lectura con algo de inercia, para no temblar.

\index{estroboscopio}%
El estroboscopio es el más sensible, y es uno de verdad. Un estroboscopio
mecánico gira a la frecuencia de la nota buscada y lo ilumina el sonido:
afinada, su dibujo se queda quieto. Cada banda del de py-tuner hace lo mismo
con números: sigue la fase de una parcial de la cuerda y avanza una franja
por segundo por cada cent que la parcial se aparta, a la derecha si está
alta, a la izquierda si está baja. Así, una desviación de una décima de
cent, que la aguja apenas muestra, sigue moviendo despacio una banda. Una
banda se apaga cuando su parcial se pierde en el ruido. La idea viene del
estroboscopio de Entropy Piano Tuner (Haye Hinrichsen y Christoph Wick) y
de Strobie (Davorin Šego).

\index{estroboscopio!vernier}\index{estroboscopio!armónicos}%
Tiene dos modos, que se eligen en la configuración o con \key{m} en la
terminal:

\begin{description}
\item[Vernier] (el de siempre) las tres bandas siguen la parcial más clara,
  la de arriba a su propia velocidad, la del medio al doble y la de abajo al
  cuádruple. Para afinar, se gira la clavija hasta que la banda de abajo se
  quede quieta.
\item[Armónicos] cada banda sigue una parcial propia: la fundamental, la
  segunda y la tercera. Una cuerda real es algo rígida, y sus parciales
  altas suenan un poco agudas; py-tuner busca cada una donde la pone la
  rigidez de la cuerda, así que en una cuerda afinada las tres se quedan
  quietas.
\end{description}

\noindent Cuál de los dos, aguja o estroboscopio, se muestra se elige en la
configuración, o con \key{v} en la terminal.

\shot{strobe_es}{0.75}{El estroboscopio solo}

\section{La historia}
\label{sec:history}

\index{historia}\index{traza}%
Bajo el estroboscopio, la historia muestra la desviación en los últimos
diez segundos, la más reciente a la derecha, sobre la banda verde de la
tolerancia: una cuerda que se asienta tras girar la clavija, o una clavija
girada al revés, se ve enseguida. Un hueco es un momento sin nota. Se
oculta en la configuración, con \key{h} en la terminal o con
\texttt{-{}-no-trace}.

\section{Desviaciones, transposición y cejilla}
\label{sec:offsets}

\index{desviaciones}%
Algunos músicos afinan unas pocas notas un poco fuera del temperamento
igual, para que los acordes que más tocan suenen más dulces.
\menu{Configuración\then Desviaciones de las notas\elipsis} da a cada nombre
de nota su desviación, de $-50$ a $+50$ cents, en todas las octavas; la
aguja, el estroboscopio y el tono apuntan entonces a la nota así movida. En
la línea de órdenes: \texttt{-{}-offset B=-2} (repetida para cada nota;
\texttt{B=0} la quita).

\shot{offsets_es}{0.35}{Desviaciones de las notas: Si 2 cents bajo, Sol 1,5 alto, Mi 1 bajo}

\index{transposición}%
\menu{Transponer} muestra los nombres movidos unos semitonos, para un
instrumento transpositor: 2 para uno en si$\flat$, cuyo Do escrito suena
si$\flat$. Solo cambian los nombres; la altura que se mide es la misma.
\index{cejilla}%
\menu{Cejilla} es el traste de una cejilla: cada cuerda sube otros tantos
semitonos, y los botones de las cuerdas muestran las notas que entonces
suenan. En la línea de órdenes, \texttt{-{}-transpose} y \texttt{-{}-capo};
en la terminal, \key{c} mueve la cejilla.

\section{El tono de referencia}
\label{sec:tone}

\index{tono de referencia}%
Para afinar de oído, py-tuner puede tocar la nota por el sonido del
computador, como lo hace pianito (Axel Stevens): una cuerda pulsada de unos
segundos, a la frecuencia exacta que se busca, con el La de referencia, la
desviación de la nota y la cejilla. \menu{Tono}, en la ventana, toca la
cuerda fija, o si no la nota que se muestra, o si no el La 4; \key{p} hace
lo mismo en la terminal. En la línea de órdenes:

\begin{quote}
\texttt{py-tuner tone E2}\\
\texttt{py-tuner tone -i guitar 1}\quad(la primera cuerda)\\
\texttt{py-tuner tone -a 442}\quad(el La 4 a 442~Hz)
\end{quote}

\section{Configuración}
\label{sec:settings}

\index{configuración}%
\menu{Configuración\elipsis} abre los ajustes menos frecuentes:

\shot{settings_es}{0.6}{La configuración}

\begin{description}
\item[Dispositivo de entrada] la entrada de sonido que se escucha, o la
  predeterminada del sistema.
\item[Canal] la mezcla de todos los canales, o uno solo (para una interfaz
  con el instrumento en una entrada).
\item[Nombres de las notas] inglés, latino o alemán.
\item[Alteraciones] sostenidos, bemoles o ambos (F$\sharp$/G$\flat$).
\item[Mostrar] la aguja y el estroboscopio, o uno de los dos.
\item[Idioma] el idioma de la interfaz, o el del sistema; se usa desde la
  próxima vez que se abra el programa.
\item[Afinada a menos de] la tolerancia, en cents (1 si no se cambia).
\item[Ignorar por debajo de] los sonidos más débiles que esto, en dBFS, se
  ignoran ($-60$ si no se cambia). Súbalo en un cuarto ruidoso.
\item[Suavizado] cuántos segundos promedia la lectura (0,5 si no se
  cambia): más es más estable, menos sigue más rápido.
\item[Cuadro de análisis] la longitud de cada cuadro, en períodos de la
  nota (16 si no se cambia): más es más preciso, menos reacciona más rápido.
\item[Estroboscopio] vernier o armónicos (sección~\ref{sec:strobe}).
\item[Mostrar la historia] la historia de la desviación
  (sección~\ref{sec:history}).
\item[Transponer, Cejilla, Desviaciones] sección~\ref{sec:offsets}.
\end{description}

\noindent Cada cambio se guarda en el acto, y la ventana, la terminal y la
línea de órdenes comparten la misma configuración.

\section{La terminal}
\label{sec:tui}

\index{terminal}\index{tui@TUI}%
\texttt{py-tuner tui} muestra, en toda la terminal, lo mismo que la
ventana: la nota en letras grandes, la frecuencia y la desviación, una
aguja de texto, el estroboscopio, las cuerdas y el nivel.

\shot{tui_es}{0.9}{La interfaz de terminal, con la cuerda Sol de una guitarra 0,4 cents alta}

\noindent Se maneja con estas teclas, y cada cambio se guarda en la
configuración:

\begin{longtable}{@{}lp{0.7\textwidth}@{}}
\toprule
Tecla & Lo que hace\\
\midrule
\endhead
\key{q} & Salir \\
\key{← / →} & Bajar / subir el La de referencia (A4) 0,1 Hz \\
\key{↓ / ↑} & Bajar / subir el La de referencia (A4) 1 Hz \\
\key{0} & Volver a poner el La de referencia (A4) en 440 Hz \\
\key{i} & El instrumento siguiente \\
\key{t} & La afinación siguiente del instrumento \\
\key{s} & Fijar la cuerda siguiente (tras la última, vuelve a la cuerda más cercana) \\
\key{n} & El sistema siguiente de nombres de notas: inglés, latino, alemán \\
\key{b} & Sostenidos, bemoles o ambos \\
\key{v} & Aguja y estroboscopio, la aguja sola o el estroboscopio solo \\
\key{m} & El estroboscopio: vernier (una parcial, tres velocidades) o armónicos (parciales 1, 2, 3) \\
\key{h} & Mostrar u ocultar la historia de la desviación \\
\key{c} & Subir la cejilla un traste (tras el 7.º, se quita) \\
\key{p} & Tocar el tono de referencia: la cuerda fija, o la nota que se muestra, o el La 4 \\
\key{l} & El idioma siguiente de la interfaz \\
\bottomrule
\end{longtable}

\section{La línea de órdenes}
\label{sec:cli}

\index{línea de órdenes}\index{cli@CLI}%
\texttt{py-tuner cli} escribe la lectura en una línea que se reescribe
(o, en una tubería, una línea cada cuarto de segundo):

\begin{verbatim}
A2       109,52 Hz   -7,5 cents  [---------*|----------]  cuerda 2
\end{verbatim}

\noindent \texttt{-{}-once} espera la primera lectura estable, la escribe y
termina, para usarla en guiones.

\shot{cli_es}{0.95}{La línea de órdenes: una grabación medida nota por nota, y la primera lectura estable}

\index{grabaciones}\index{wav@WAV}%
Con un archivo WAV, \texttt{py-tuner cli -{}-file ARCHIVO.wav} analiza toda
la grabación de una vez y escribe una línea por cada nota que suena en
ella: cuándo suena, la nota, su frecuencia, la que debería tener, la
desviación y su dispersión:

{\footnotesize
\begin{verbatim}
0,20-   1,77 s  E2  82,550 Hz  (objetivo 82,407)  +3,00 cents  ±0,02  cuerda 1
1,80-   3,52 s  A2  109,52 Hz  (objetivo 110,00)  -7,50 cents  ±0,03  cuerda 2
\end{verbatim}
}

\noindent La primera décima de segundo de cada nota se deja fuera, porque
una cuerda pulsada suena algo alta al principio. Con
\texttt{-{}-realtime}, en cambio, el archivo suena a su ritmo, como si
viniera de la tarjeta de sonido; \texttt{-{}-file} funciona igual con la
ventana y con la terminal.

Las opciones comunes a las tres maneras:

\begin{longtable}{@{}lp{0.62\textwidth}@{}}
\toprule
Opción & Lo que hace\\
\midrule
\endhead
\texttt{-a}, \texttt{-{}-a4 HZ} & el La de referencia (A4), de 380 a 480 Hz\\
\texttt{-i}, \texttt{-{}-instrument} & \texttt{chromatic} (si no se cambia), \texttt{guitar}, \texttt{bass}, \texttt{double-bass} o \texttt{violin}\\
\texttt{-t}, \texttt{-{}-tuning NOMBRE} & la afinación, por su nombre (p.~ej.\ \texttt{"Drop D"})\\
\texttt{-s}, \texttt{-{}-string N} & comparar solo con la cuerda N (1 = la más grave); \texttt{auto}, con la más cercana\\
\texttt{-{}-names} & \texttt{english}, \texttt{latin} o \texttt{german}\\
\texttt{-{}-accidentals} & \texttt{sharps}, \texttt{flats} o \texttt{both}: sostenidos, bemoles o ambos (F\#/Gb); también \texttt{-{}-flats}, \texttt{-{}-sharps}\\
\texttt{-{}-tolerance CENTS} & afinada a menos de estos cents\\
\texttt{-d}, \texttt{-{}-device DISP} & la entrada, por su número (\texttt{py-tuner devices}) o parte de su nombre\\
\texttt{-r}, \texttt{-{}-rate HZ} & la frecuencia de muestreo (la del dispositivo, si no se da)\\
\texttt{-c}, \texttt{-{}-channel} & \texttt{mix}, o un número de canal desde 1\\
\texttt{-f}, \texttt{-{}-file WAV} & un archivo WAV en lugar de la tarjeta de sonido\\
\texttt{-{}-min-level DBFS} & ignorar los sonidos más débiles que esto\\
\texttt{-{}-smoothing SEGUNDOS} & los segundos que promedia la lectura\\
\texttt{-{}-periods N} & los períodos de la nota en cada cuadro de análisis\\
\texttt{-{}-display} & \texttt{both}, \texttt{needle} o \texttt{strobe}\\
\texttt{-{}-offset NOTA=CENTS} & afinar un nombre de nota fuera del temperamento igual, en todas las octavas\\
\texttt{-{}-transpose N} & los nombres para un instrumento transpositor\\
\texttt{-{}-capo TRASTE} & una cejilla, que sube todas las cuerdas\\
\texttt{-{}-strobe} & \texttt{vernier} o \texttt{harmonics}\\
\texttt{-{}-trace}, \texttt{-{}-no-trace} & la historia de la desviación\\
\texttt{-{}-language} & \texttt{es} o \texttt{en}\\
\texttt{-{}-save} & guardar las opciones dadas como la nueva configuración\\
\bottomrule
\end{longtable}

\section{Cómo mide}
\label{sec:how}

\index{precisión}%
Cada análisis tiene dos pasos. Uno grueso halla el período del sonido con
el método de McLeod, que no puede errar la octava; uno fino ubica después
cada armónico de la cuerda, hasta el decimosexto, en el máximo exacto de su
espectro, y los ajusta todos a la vez al modelo de una cuerda real, algo
rígida, $f_n = n f_0 \sqrt{1 + B n^2}$. Los armónicos fuertes cuentan más
que los débiles, y los que se hunden en el ruido quedan fuera. La lectura
que se muestra es la media recortada del último medio segundo de cuadros.

\index{octava}%
Una cuerda real puede engañar a un período: dos cuerdas que resuenan
comparten un período una octava o una duodécima bajo la nota, y una cuerda
de acero cuya fundamental el micrófono apenas recoge se repite casi cada
medio período. Por eso los candidatos alrededor del período (un tercio, la
mitad, el doble y el triple) se puntúan por lo que sobresalen sus armónicos
del ruido que los rodea, menos lo que hay donde caerían los armónicos de
una octava equivocada, una idea tomada de Strobie y de Entropy Piano Tuner;
la nota que ya suena lleva ventaja mientras se apaga, como en Tuneo
(Braulio Ríos). Una nota nueva justo después de pulsar se toma enseguida;
sin pulsación, debe durar 0,3 segundos, o 0,6 si es la octava, la duodécima
o la quinta de la que suena. Y con las cuerdas de un instrumento solo se
buscan las notas cercanas a ellas. En once grabaciones de cuerdas reales de
guitarra, estas reglas dejaron una nota falsa donde había quince.

\index{consejos}%
Algunas costumbres ayudan:

\begin{itemize}
\item Tocar la cuerda sola, y dejar pasar el primer instante del pulsado:
  una cuerda pulsada suena algo alta al principio, y luego se asienta.
\item Acercar el micrófono al instrumento, lejos de ventiladores y otros
  ruidos; si el cuarto es ruidoso, subir \menu{Ignorar por debajo de}.
\item Con las cuerdas graves, darle a la lectura medio segundo más: el
  cuadro de un Mi grave del bajo dura casi 0,4~s.
\item Para afinar fino, mirar la banda de abajo del estroboscopio.
\end{itemize}

\section{Los archivos que guarda}
\label{sec:files}

\index{archivos}%
El programa guarda un solo archivo:
\texttt{\textasciitilde/.config/py-tuner/settings.json} (bajo
\texttt{\$XDG\_CONFIG\_HOME}, si está definida), con toda la configuración,
en JSON. Se puede editar a mano con el programa cerrado; un valor que no
tenga sentido se cambia por el de fábrica, así que el programa siempre
arranca. Borrarlo devuelve toda la configuración de fábrica.

\section{Cuando algo sale mal}

\index{problemas}%
\begin{description}
\item[No se oye nada (``Toque una nota'')] la barra de nivel no se mueve:
  la entrada elegida no es la del instrumento. Elija otra en
  \menu{Configuración} (\texttt{py-tuner devices} las lista), y revise que
  la entrada no esté silenciada en la configuración de sonido del sistema.
\item[El nivel se mueve, pero no aparece ninguna nota] el sonido es más
  débil que \menu{Ignorar por debajo de}, o demasiado ruidoso para tener
  altura. Acerque el micrófono, o baje \menu{Ignorar por debajo de}.
\item[La nota salta a la cuerda siguiente] la cuerda está tan floja que
  queda más cerca de otra; fije la cuerda que está afinando
  (sección~\ref{sec:window}).
\item[La lectura difiere unos cents de la de otro afinador] revise que los
  dos tengan el mismo La de referencia.
\item[``No se puede abrir la entrada''] otro programa tiene la tarjeta de
  sonido, o el dispositivo se desconectó; el mensaje dice la causa.
\item[El tono no suena] la salida de sonido del computador está
  silenciada, u otro programa la tiene; el mensaje dice la causa.
\item[Un ajuste está mal] borre el archivo de configuración
  (sección~\ref{sec:files}).
\end{description}

\noindent
Los problemas y los deseos pueden escribirse como \emph{issues} en GitHub o
en Codeberg (las direcciones están en la primera página), o enviarse a
\href{mailto:nicolasvaughan@proton.me}{\texttt{nicolasvaughan@proton.me}}.

\printindex

\end{document}

% ===============================================
%%% Local Variables:
%%% mode: latex
%%% TeX-master: t
%%% End:
